\section{ОСНОВНАЯ ЧАСТЬ}



\subsection{Утилита для сбора логов по camunda events}

\subsubsection{Задача}
\begin{itemize}
  \item Сервис реализует API для получения истории событий Camunda за произвольный временной диапазон с параметрами startTime и endTime, передаваемыми через URI.
  \item Возвращаемые данные включают только активности типов: service task, user task, exclusive gateway, send task, start event, script task, inclusive gateway, intermediate message throw event.
  \item Формат выгрузки данных — .xlsx, со столбцами Activity ID, Activity Name, Start Time, End Time, Duration (ms), Activity Type.
  \item Сбор истории выполняется без заметного влияния на производительность основной системы (разница в средней загрузке CPU не более 5\%).
  \item Ошибки при запросе возвращают корректный HTTP-ответ с кодами: 400 (некорректные параметры), 500 (внутренняя ошибка).

\end{itemize}

\subsubsection{Аналитика задачи}

Задача была получена и начата на одном из обсуждений текущих приоритетов команды 16.02.2026.
Даннуя задачу начали обсуждать в предверии нагрузочного тестирования всей системы, для сверки метрик,
которые были поставлены в документации проекта. Так как аудит должен был включать проверку каждого сервиса системы, 
то требовалось проверить, не является ли система бизнес процессов узким местом системы. После получения задачи, я связался с аналитиком, который 
мне более детально объяснил, что требуется от механизма сбора данных, а если конкретнее, то нужно было собрать все 
активности для отслеживания токена процесса следующего вида: service task, user task, exclusive gateway, send task, start event,
script task, inclusive gateway, intermediate message throw event, так же требовалось дополнительная информация о номере процесса, его имени, времени начала, 
времени конца, затраченного времени на выполенения конкретной активности и её тип. 
Получение этой информации аналитик попросил реализовать интерфейс ввиде простых запросов по протоколу передачи гипертекста. Для получения данных был предложен следующий формат запроса:

\begin{lstlisting}[
  caption={Запрос лога активности},
  label={lst:activity-log}
]
exmaple.ru/api/activity-log?startTime=2026-06-11T11:00:00&endTime=2026-06-11T14:28:00
\end{lstlisting}

Данный запрос включает в себя параметры, который решено передавать через URI. 
Параметры включают: время с которого стоит делать сбор активностей и время окончания сбора данных.
Указание временного промежутка должно помочь аналитикам при работе с несколькими процессами, так как кол-во строк истории часто измеряется в сотнях тысяч. 
Так же это позволит делать срезы, наиболее интересные для исследования в точках зависания, чтобы просмотреть, какие конкретно активности исполнялись.
Формат для представления данных был выбран ввиде exel-отчёта.

\subsubsection{Заведение задачи в системе трекинга}
После перевода задачи аналитики в статус "Выполнена". 
Я ещё раз ознакомился с аналитикой и создал задачу для разработки, прикрепил в "Относится к" задачу аналитики. 
Проставил срок для выполнения с 16.02.2026 по 05.03.2026, в этот срок так же закладываю время на тестирование нового механизма. 
Задаче были назначены следующие метки -- бизнес-логика, задача и проставлена веха, в которую попадёт задача. 
Критичность выставлена "Средний", так как задача не блокирует другие задачи и не является багом.
После описания задачи, я перевёл её в статус "В процессе".

\subsubsection{Исследование текущего интерфейса системы бизнес процессов}
Для реализации задачи нужно изначально изучить интерфейс, 
который предоставляет нам система бизнес процессов для запроса истории. На официальном сайте с описанием интерфейсов [1] я нашёл нужный мне метод для сбора истории. 
Данный метод поддерживает фильтрацию по временному промежутку с помощью параметров запроса, 
что полностью соответствует требованию аналитика о возможности сбора активностей за определённый временной интервал. 
Кроме того, метод возвращает следующие поля для каждой активности:
\begin{enumerate}
  \item id — уникальный идентификатор экземпляра активности;
  \item activityId — идентификатор активности;
  \item activityName — имя активности;
  \item activityType — тип активности;
  \item processDefinitionId — идентификатор определения процесса;
  \item processInstanceId — номер экземпляра процесса;
  \item startTime — время начала выполнения активности;
  \item endTime — время окончания выполнения активности;
  \item duration — затраченное время в миллисекундах.
\end{enumerate}
Данные поля нам подходят по бизнес требованиям. Так как мы используем клиентскую библиотеку для взаимодействия с отдельными модулем камунды, то имеет смысл воспользоваться инверсией контроля.
Сам сервис мы будет использовать при помощи внедрения зависимостей через конструктор при помощи аннотирования класса для снижения прогретого кода. У клиента есть несколько методов, которые подходят нам, но я выбираю метод, где на вход возможны 3 параметра: номер страницы для пагинации, размер страницы и объект с параметрами фильтрации запроса.
Данный метод позволяет нам сразу передать в теле запроса объект с фильтрами, включая параметры $startedAfter$ и $startedBefore$. Возвращает он список объектов, содержащих все необходимые поля: идентификатор активности, её тип, время начала и окончания, длительность выполнения в миллисекундах, а также идентификаторы экземпляра процесса и определения процесса. 
Всё это даёт возможность собирать активности только за указанный аналитиком период времени, что было одним из ключевых требований задачи.
Выбранный метод использует POST-запрос. Параметры временного диапазона передаются в теле запроса, а не в URL.

\subsubsection{Контроллер}

Реализацию задачи я решил начать сверху вниз, начиная с интерфейса у бизнес логики. 
На проекте используются стандартные для фреймворка контроллеры, запрос в которые приходит из специального сервлета, который находит конкретный конроллер и метод по заданному пути.
Разработка контроллера началась с определения точки входа. Сам тип запроса выбран как $GET$, так как операция не изменяет данные на уровне сервиса, а только возвращает их, это соответствует семантике типа метода.
Для поддержки временной фильтрации в контроллер были добавлены необязательные параметры: время начала и время конца. 
Они принимаются в формате в стандартном ISO-формате, как это было заданно в аналитической задаче. 
Оба параметра помечены как необязательные, чтобы при необходимости можно было выгрузить полный журнал без временных ограничений.
Добавил валидацию для ошибочных запросов, которая учитывает:
\begin{enumerate}
  \item начальное время не может быть позже конечного;
  \item начальное время не может быть в будущем;
  \item конечное время не может быть в будущем.
\end{enumerate}
При нарушении любого из этих условий будет ответ с кодом 400, сообщая клиенту о неверно сформированном запросе.
В контроллер добавил логирование для отслеживания факта обращения к эндпоинту, а также логирование ошибок и предупреждений. 
Это позволяет в дальнейшем анализировать работу механизма.
Контроллер не занимается непосредственно сбором данных, поэтому я решил делегировать эту работу сервису, вызывая его метод, так как это соответствует принципу разделения ответственности.
После получения файла от сервиса контроллер проверяет, не равен ли результат пустому значению. 
Если данных за указанный период не найдено, возвращается ответ с кодом 404. 
Если файл успешно создан, контроллер формирует имя файла с текущей датой и временем, 
чтобы каждый скачанный файл имел уникальное имя и не перезаписывал предыдущие.
Для возврата файла контроллер использует потоковый режим, который позволяет не загружать его полностью в память. 
Это важно при работе с большими объёмами данных. В заголовки ответа добавляется тип для Excel-файлов. Код ответа — 200.
В случае возникновения ошибки при генерации файла (например, проблем с доступом к Camunda) контроллер возвращает ответ с кодом 500, 
логируя при этом детали ошибки для последующего анализа.
По требованиям ревью, я сделал документацию контроллеру с помощью аннотаций Swagger. 
В описании указаны возможные коды ответов, форматы параметров и описание работы метода.
Сам код контроллера прикреплён в \textbf{Приложении А}.

\subsubsection{Сервис}

Для реализации самой выгрузки данных и формировании отчёта, 
был разработан сервис. Его основное назначение -- собрать данные об активностях за заданный промежуток времени и 
сформировать exel-файл с этими данными. Так же дополнительно из конфигурации были вытянуты человекочитаемые совпадения из бизнес процесса для сопоставления с 
уникальными названиями активностей. При запуске сервиса, модуль подгружает человекочитаемые названия из файла, который лежит в classpath. 
Загрузка выполняется однократно при помощи введения некоторой аннотации $PostConstruct$. 
Само местоположение файла можно задавать через конфигурационный параметр для гибкости использования.
Предзагрузка сделана для создания закешированного вида файла, чтобы ускорить работу с сопоставлением имён.
Для получения исторических записей, я использовал клиентскую библиотеку, внедрение зависимостеи я сделал посредством аннотации $RequiredArgsConstructor$.
Я сделал создание транспортного объекта для запроса, куда включил фильтры, которые были переданы: $startedAfter$ и $startedBefore$.
Если параметры не переданы, фильтры не применяются, и выгружаются все доступные активности.
Из полученного от Camunda списка активностей отфильтровываются только те, у которых заполнено поле $endTime$. 
Я сделал такую фильтрацию так как незавершённые активности не имеют полной информации о времени выполнения и длительности, что делает их бесполезными для анализа.
Так же я предусмотрел излишнюю работу и если список пустой, то пользователю вернётся информативный ответ с ошибкой, что за заданный промежуток ничего не найдено.
Для формирования выходного файла я использовал библиотекту $Apache POI$, которая позволила мне генерировать файлы нужного формата.
Сервис создаёт новую книгу, добаляет туда лист, формирует строку заголовков и далее заполняет таблицу.
При заполнении каждой строки я выполняю некоторые преобразования: 
\begin{enumerate}
  \item Даты startTime и endTime форматируются в человекочитаемый вид.
  \item Технические идентификаторы активностей преобразуются в понятные названия, который обращается к загруженному при старте справочнику.
  \item Для родительских активностей также извлекаются их названия и схемы.
\end{enumerate}
При формировании файла они сохраняются во временную директорию, это я сделал чтобы
не оставлять файлы на диске после завершения запроса пользователя. При возникновении ошибок сервис должен вывести тело ошибки и выбрасывать исключение, которое далее обработается на уровне контроллера.
После записи файла книга закрывается, а поток вывода завершается с помощью конструкции $ try-with-resources$, что так же будет гарантировать освобождение системных ресурсов.

Файл с человекочитаемыми названиями имеет следующую структуру:

\begin{lstlisting}[
  caption={Файл с человекочитаемыми названиями},
  label={lst:activity-log}
]
{
  "ServiceTask_1": {
    "name": "Проверка данных",
    "schema": "Проверка входных параметров"
  },
  "UserTask_1": {
    "name": "Согласование",
    "schema": "Утверждение заявки руководителем"
  }
}
\end{lstlisting}
Сам код сервиса можно увидет в \textbf{Приложении А}.

\subsubsection{Затронутое API}

В ходе реализации я задействовал следующие компоненты REST API Camunda:
\begin{itemize}
\item POST $/history/activity-instance$ — метод, выполняющий запрос исторических записей активностей.
В теле запроса я передаю параметры фильтрации:
$startedAfter$ и $startedBefore$.

\item Клиентская библиотека Camunda 7 Java Client — предоставляет обёртку $HistoricActivityInstanceApi$,
которая инкапсулирует формирование HTTP-запроса и десериализацию ответа.
\end{itemize}

\subsubsection{Допустимые ошибки}

В процессе сбора логов возможны некоторые ошибки, 
обработку которых я предусморел и представил в Таблице 1.1. 

\begin{table}[H]
  \centering
  \caption{Коды ответов и их значения}
  \begin{tabular}{|p{3cm}|p{10cm}|}
    \hline
    \textbf{Код} & \textbf{Ситуация} \\
    \hline
    400 & $startTime$ позже $endTime$; любая из дат находится в будущем; неверный формат даты. \\
    \hline
    404 & За указанный период не найдено ни одной завершённой активности (список пуст после фильтрации). \\
    \hline
    500 & Ошибка подключения к Camunda; ошибка чтения JSON-файла с маппингом имён; ошибка записи временного Excel-файла; необработанное исключение API. \\
    \hline
  \end{tabular}
\end{table}

\subsubsection{Замер потребления ресурсов при нагрузке}

Замер я проводил на специальном контуре, где были созданы процессы. Один из процессов я запустил вручную, 
чтобы начала щаписываться история. Сам процесс был запущен в 22:49 (время специфично для сервера и не совпадает с реальным $GMT+3$).
После завершением процесса в 22:53, я сделал запрос на сбор истории в 22:54.
\begin{lstlisting}[
  caption={Запрос лога активности},
  label={lst:activity-log}
]
gpz.prod.ru/api/activity-log?startTime=2026-02-22T15:00:00&endTime=2026-02-27T18:00:00
\end{lstlisting}

Для рассмотрения нагрузки я воспользовался готовым инструментом сбора метрик по $CPU Usage$ и $Memory Usage$.

\begin{figure}[H]
  \centering
  \includegraphics[
    width=\linewidth,
    height=0.9\textheight,
    keepaspectratio
  ]{\detokenize{artifacts/back-cpu.png}}
  \caption{Скриншот графика загрузки процессора сервиса со сбором активностей}
\end{figure}

\begin{figure}[H]
  \centering
  \includegraphics[
    width=\linewidth,
    height=0.9\textheight,
    keepaspectratio
  ]{\detokenize{artifacts/back-mem.png}}
  \caption{Скриншот графика загрузки памяти сервиса со сбором активностей}
\end{figure}

\begin{figure}[H]
  \centering
  \includegraphics[
    width=\linewidth,
    height=0.9\textheight,
    keepaspectratio
  ]{\detokenize{artifacts/camunda-cpu.png}}
  \caption{Скриншот графика загрузки процессора сервиса с camunda}
\end{figure}

\begin{figure}[H]
  \centering
  \includegraphics[
    width=\linewidth,
    height=0.9\textheight,
    keepaspectratio
  ]{\detokenize{artifacts/camunda-mem.png}}
  \caption{Скриншот графика загрузки памяти сервиса с camunda}
\end{figure}

Данная утилита не оказывает значимое влияние на загрузку системы,
поэтому её можно безопасно использовать для сбора и аналитики данных.

\subsubsection{Использованные технологии}

При реализации были использованы следующие технологии и библиотеки:

\begin{enumerate}
  \item Java 17 — язык реализации.
  \item Spring Boot 2.7 — фреймворк для создания REST-контроллера и внедрения зависимостей.
  \item Camunda 7 Java Client — для вызова исторического API.
  \item Apache POI (XSSFWorkbook) — для генерации Excel-файлов.
  \item Jackson — для парсинга JSON-файла с человекочитаемыми именами активностей.
  \item Swagger — для документирования эндпоинта.
  \item SLF4J — для логирования.
\end{enumerate}

\subsubsection{Почему были выбраны технологии}

Некоторые технологии уже использовались на проекте, поэтому было решено взять из для избежания добавления новых зависимостей в проект. Технологии которые уже были использованы -- протестированы на продуктивных стендах.

\begin{itemize}
  \item Apache POI выбран, так как он является стандартом для работы с Excel в Java, поддерживает потоковую запись больших файлов.

  \item Camunda 7 Java Client использовался вместо прямых HTTP-вызовов, так как клиент поддерживается официальным сообществом инструмента.

  \item Jackson взят как стандартный инструмент Spring Boot — не требует дополнительных зависимостей.

  \item Swagger — является достаточно легковестным расширением OpenAPI для документирования кода.
\end{itemize}

\subsubsection{Вывод}

Сделан сервис, который реализует API для получения истории событий Camunda за произвольный временной диапазон с параметрами startTime и endTime, передаваемыми через URI. Согласно требованиям -- возвращаемые данные включают активности типов: service task, user task, exclusive gateway, send task, start event, script task, inclusive gateway, intermediate message throw event.
Сделан запрашиваемый формат выгрузки — .xlsx, со столбцами Activity ID, Activity Name, Start Time, End Time, Duration (ms), Activity Type.
После разработки был проведён замер, который показал, что сбор истории выполняется без заметного влияния на производительность основной системы и не загружает процессор более чем на 5%.
Ошибки при запросе возвращают корректный HTTP-ответ с кодами: 400 (некорректные параметры), 500 (внутренняя ошибка).


% Изначально задача была затрона на одном из собраний команды 16.02.2026. 
% Где происходил опрос разработчиков по прогрессу их задачек, далее предлагались задачи на выполнение, 
% но так как задача входит в мои задачи практики, то мне её выдали на разработку.
% Сама задача изначально вызникла из проблемы, в одном из сервисов был метод, 
% который каждые 0.5 секунды делал опрос истории камунды

% поменяй название некоторых переменных, чтобы этот код был обезличен
% @Scope(value = ConfigurableBeanFactory.SCOPE_SINGLETON)
% @Service
% @RequiredArgsConstructor
% public class ScheduledBpmnHistoryServiceBean {


%   private final AtomicBoolean isRunning = new AtomicBoolean();


%   private final WorkflowStepInstanceService workflowStepInstanceFacade;


%   private final WorkflowObjectService workflowObjectFacade;


%   private final CamundaAvailable camundaAvailable;
%   @Lazy
%   private final CamundaTaskFetchUtilServiceBean camundaTaskFetchUtilService;


%   @Scheduled(fixedDelay = 500)
%   //@NuxeoClientThreadAuthorization
%   protected void readHistory() {
%     if (Boolean.TRUE.equals(camundaAvailable.isActive()) && isRunning.compareAndSet(false, true)) {
%       try {
%         Date latestRequestedDate = workflowStepInstanceFacade.latestRequestedDate();
%         latestRequestedDate = camundaTaskFetchUtilService.fixStartedAfter(latestRequestedDate);
%         camundaTaskFetchUtilService.processHistoricActivityInstances(latestRequestedDate,
%             null);
%         camundaTaskFetchUtilService.processUpcomingTaskInstances(latestRequestedDate, null);
%         Date latestEndDate = workflowObjectFacade.latestEndDate();


%         while (camundaTaskFetchUtilService.processHistoricProcessInstances(latestEndDate)) {
%           latestEndDate = workflowObjectFacade.latestEndDate();
%         }
%       } finally {
%         isRunning.compareAndSet(true, false);
%       }
%     }
%   }
% }


%   public Date fixStartedAfter(Date startedAfter) {
%     lastDiffUpdate.ifPresentOrElse(lastDiff -> {
%       this.startedAfterDiff.setTime(
%           startedAfter.getTime() - (startedAfter.getTime() - lastDiff.getTime()));
%       this.lastDiffUpdate = Optional.of(startedAfter);
%     }, () -> {
%       this.lastDiffUpdate = Optional.ofNullable(startedAfter);
%       this.startedAfterDiff = startedAfter;
%     });
%     return startedAfterDiff;
%   }

% public Date latestRequestedDate() {
%     var w = latestByAttribute(WorkflowStepInstance.REQUESTED_DATE);
%     return w != null ? w.getRequestedDate().orElse(null) : null;
%   }

% public void processHistoricActivityInstances(final Date startedAfter,
%       final Date startedBefore) {
%     HistoricActivityInstanceQueryDto query = new HistoricActivityInstanceQueryDto();
%     Optional.ofNullable(startedAfter).ifPresent(query::setFinishedAfter);
%     Optional.ofNullable(startedBefore).ifPresent(query::setFinishedBefore);


%     Map<WorkflowObject, String> history = new HashMap<>();
%     try {
%       var intermediateHistory = historicActivityRestService.queryHistoricActivityInstances(
%               null, null, query).stream()
%           .filter(Objects::nonNull)
%           .filter(h -> h.getEndTime() != null)
%           .map(HistoricActivityInstanceDto::getProcessInstanceId).collect(Collectors.toSet())
%           .stream().map(id -> {
%             try {
%               return processInstanceRestService.getActivityInstanceTree(id);
%             } catch (ApiException e) {
%               return null;
%             }
%           })
%           .filter(Objects::nonNull)
%           .map(ActivityInstanceDto::getChildActivityInstances).filter(a -> !a.isEmpty())
%           .map(a -> a.get(0)).filter(this::filter).map(a -> Pair.of(a,
%               workflowObjectFacade.findBaseByProcessInstanceId(a.getProcessInstanceId())))
%           .filter(p -> p.getB().isPresent())
%           .map(p -> workflowObjectFacade.createWorkflowObjectFromNuxeo(p.getB().get())
%               .thenApply(f -> Pair.of(p.getA(), f)))
%           .toList();
%       try {
%         CompletableFuture.allOf(intermediateHistory.toArray(CompletableFuture[]::new)).get();
%       } catch (InterruptedException e2) {
%         LOGGER.error("", e2);
%         Thread.currentThread().interrupt();
%       } catch (ExecutionException | CompletionException | CancellationException e) {
%         LOGGER.error("", e);
%       }
%       history = intermediateHistory.stream().map(CompletableFuture::join)
%           .filter(p -> !Objects.equals(p.getA().getActivityId(), p.getB().getWorkflowStatus()))
%           .collect(Collectors.toMap(Pair::getB,
%               p -> getPrefix(p.getA().getActivityType()) + convertId(p.getA().getActivityId())))
%           .entrySet().stream().filter(e -> !e.getKey().getWorkflowStatus().equals(e.getValue()))
%           .collect(Collectors.toMap(Map.Entry::getKey, Map.Entry::getValue));
%     } catch (ApiException e) {
%       LOGGER.error(e.getResponseBody(), e);
%     }


%     history.forEach((workflowObject, state) -> {
%       var changed = new AutoHolder<>(workflowObject::getInput);
%       changed.getValue().setWorkflowTransition(state);
%       if (changed.isPresent()) {
%         workflowObjectSubscriberBean.workflowPending(
%             new WorkflowPending(workflowObject, state, ""));
%       }
%     });
%   }

% public void processUpcomingTaskInstances(final Date startedAfter, final Date startedBefore) {
%     processPendingSteps();


%     HistoricTaskInstanceQueryDto query = new HistoricTaskInstanceQueryDto();
%     query.setUnfinished(true);
%     query.setUnassigned(true);
%     Optional.ofNullable(startedAfter).ifPresent(query::setStartedAfter);
%     Optional.ofNullable(startedBefore).ifPresent(query::setStartedBefore);


%     List<HistoricTaskInstanceDto> historicTaskInstances;
%     try {
%       historicTaskInstances = historicTaskInstanceRestService.queryHistoricTaskInstances(
%           null, null, query);
%     } catch (ApiException e2) {
%       LOGGER.error(e2.getResponseBody(), e2);
%       return;
%     }
%     Map<String, List<HistoricTaskInstanceDto>> processInstanceId2history = historicTaskInstances
%         .stream().collect(Collectors.groupingBy(HistoricTaskInstanceDto::getProcessInstanceId));
%     Map<String, CompletableFuture<WorkflowObject>> processInstanceId2WorkflowObject =
%         workflowObjectFacade.findByProcessInstanceIds(processInstanceId2history.keySet());
%     try {
%       CompletableFuture.allOf(processInstanceId2history.entrySet().stream().map(
%           e -> Optional.ofNullable(processInstanceId2WorkflowObject.get(e.getKey()))
%               .map(f -> f.thenApply(Optional::of))
%               .orElseGet(() -> ThreadScopeHelper.thenApplyAsync(
%                   camundaTaskDataBean.getTaskVariableValues(
%                       e.getValue().stream().findAny().orElseThrow().getId())
%                       .thenCompose(v -> v.entrySet().stream()
%                           .filter(ev -> ev.getKey().endsWith("__common__id"))
%                           .map(Map.Entry::getValue).map(String.class::cast).findFirst()
%                           .map(workflowObjectFacade::findByDataObjectIdAsync)
%                           .orElseGet(() -> CompletableFuture.completedFuture(
%                               Optional.<WorkflowObject>empty()))),
%                   wo -> wo.map(
%                       w -> {
%                         LOGGER.info("candidate orphan process instance id=" + e.getKey());
%                         return w;
%                       })
%                       .or(() -> {
%                         LOGGER.error("delete orphan process instance id=" + e.getKey());
%                         // ProcessInstanceResource processInstance = processInstanceRestService.get(
%                         // ProcessInstanceRestService.PATH + "/" + e.getKey());
%                         // processInstance.deleteProcessInstance(false, false, false);
%                         try {
%                           processInstanceRestService.deleteProcessInstance(e.getKey(), false, false,
%                               false, false);
%                         } catch (ApiException e1) {
%                           LOGGER.error(e1.getResponseBody(), e1);
%                         }
%                         return Optional.empty();
%                       })))
%               .thenCompose(wo -> wo
%                   .map(workflowObject -> CompletableFuture.allOf(e.getValue().stream()
%                       .map(historicDto -> {
%                         var workflowStepTemplate = workflowObject.getWorkflowType()
%                             .getWorkflowStepTemplates().stream()
%                             .filter(wst -> wst.getTaskDefKey().equals(
%                                 historicDto.getTaskDefinitionKey()))
%                             .findAny();
%                         if (workflowStepTemplate.isPresent()
%                             && Bpmn.Lifecycle.STARTING_PROCESS.equals(
%                                 workflowObject.getWorkflowStatus())) {
%                           putToPendingSteps(historicDto.getId(),
%                               workflowObject.getWorkflowObjectID(),
%                               workflowStepTemplate.get());
%                           return Optional.<CompletableFuture<Boolean>>empty();
%                         } else {
%                           return Optional.of(bpmnWorkflowStepUtilityService.getObject()
%                               .processIncomingTask(e.getKey(), workflowObject, historicDto,
%                                   workflowStepTemplate));
%                         }
%                       })
%                       .filter(Optional::isPresent)
%                       .map(Optional::get).toArray(CompletableFuture[]::new)))
%                   .orElseGet(() -> CompletableFuture.completedFuture(null))))
%           .toArray(CompletableFuture[]::new)).get();
%     } catch (InterruptedException e2) {
%       LOGGER.error("", e2);
%       Thread.currentThread().interrupt();
%     } catch (ExecutionException | CompletionException | CancellationException e) {
%       LOGGER.error("", e);
%     }
%   }


% public boolean processHistoricProcessInstances(final Date finishedAfter) {
%     HistoricProcessInstanceQueryDto query = new HistoricProcessInstanceQueryDto();
%     query.setFinished(true);
%     query.setActive(false);
%     Optional.ofNullable(finishedAfter).ifPresent(query::setFinishedAfter);
%     List<HistoricProcessInstanceDto> historicProcessInstances;
%     try {
%       historicProcessInstances = historicProcessInstanceRestService.queryHistoricProcessInstances(
%           null, null, query);
%     } catch (ApiException e) {
%       LOGGER.error(e.getResponseBody(), e);
%       return true;
%     }
%     Map<String, HistoricProcessInstanceDto> instances = historicProcessInstances.stream()
%         .collect(Collectors.toMap(HistoricProcessInstanceDto::getRootProcessInstanceId, h -> h,
%             (a, b) -> a));
%     var resultChanged = Holder.of(false);
%     CompletableFuture.allOf(
%         workflowObjectFacade.findByProcessInstanceIds(instances.keySet()).values()
%         .stream()
%         .map(ff -> ff.thenCompose(workflowObject -> {
%           HistoricProcessInstanceDto historicDto =
%               instances.get(workflowObject.getProcessInstanceId());
%           var changed = new AutoHolder<>(workflowObject::getInput);
%           if (workflowObject.getEndDate() == null) {
%             LOGGER.info("workflowObjectID={} endDate={}", workflowObject.getWorkflowObjectID(),
%                 historicDto.getEndTime());
%             changed.getValue().setEndDate(historicDto.getEndTime());
%           }
%           var changedState = Holder.of(false);
%           switch (workflowObject.getWorkflowStatus()) {
%             case Bpmn.Lifecycle.FINISHED, Bpmn.Lifecycle.CANCELED:
%               break;
%             default:
%               // processing of PAUSED_BY_ERROR is unnecessary
%               // because process instance has been terminated/completed
%               switch (historicDto.getState()) {
%                 case EXTERNALLY_TERMINATED:
%                   changed.getValue().setWorkflowTransition(Bpmn.Lifecycle.CANCELED);
%                   break;
%                 default:
%                   LOGGER.info("workflowObjectID={} status={}", workflowObject.getWorkflowObjectID(),
%                       workflowObject.getWorkflowStatus());
%                 case COMPLETED:
%                   changed.getValue().setWorkflowTransition(Bpmn.Lifecycle.FINISHED);
%                   break;
%               }
%               changedState.setValue(true);
%               break;
%           }
%           if (changed.isPresent()) {
%             return workflowObjectFacade.findByPrimaryKey(workflowObject.getWorkflowObjectID())
%                 .thenCompose(updatedWorkflowObject -> {
%                   resultChanged.setValue(true);
%                   if (Boolean.TRUE.equals(changedState.getValue())) {
%                     return workflowObjectSubscriberBean.workflowFinished(
%                         new WorkflowFinished(updatedWorkflowObject,
%                             historicDto.getState().name()
%                         ));
%                   }
%                   return CompletableFuture.<Void>completedFuture(null);
%                 });
%           } else {
%             return CompletableFuture.<Void>completedFuture(null);
%           }
%         }))
%         .toArray(CompletableFuture[]::new))
%         .join();
%     return resultChanged.getValue();
%   }

